\documentclass[a4paper, 12pt]{article}

\usepackage[utf8]{inputenc}
\usepackage[english]{babel}
\usepackage{hyperref}
\usepackage{graphicx}
\usepackage{subfiles}
\usepackage{blindtext}
\usepackage[margin=1in, includefoot]{geometry}
\usepackage{amsmath}
\usepackage{algorithm}
\usepackage[noend]{algpseudocode}
\usepackage{blindtext}
\usepackage{enumitem}
\usepackage[framed,numbered]{sclang-prettifier}
\usepackage[T1]{fontenc}
\usepackage{xcolor}
\usepackage{etoolbox}
\usepackage{fancyhdr} %testatine e piè pagina
\pagestyle{fancy}
    \fancyhf{}
    \chead{Euclidean Algorithm for auto-generative patterns in a Supercollider application}
    \lfoot{de Bari, Albertini}
    \rfoot{{\textsc{Computer Music}}}
    \cfoot{\thepage}
    \renewcommand{\headrulewidth}{0.7pt}
\renewcommand{\footrulewidth}{0.7pt}

\newcommand{\algorithmautorefname}{Algorithm}
\definecolor{newyellow}{RGB}{228,153,0}
\definecolor{newpink}{RGB}{255,20,147}

\renewcommand{\lstlistingname}{Code}% Listing -> Algorithm
\renewcommand{\lstlistlistingname}{List of \lstlistingname s}
\newcommand\blueref[1]{{\hypersetup{linkbordercolor=blue}\autoref{#1}}}
\newcommand\yellowref[1]{{\hypersetup{linkbordercolor= newyellow}\autoref{#1}}}
\newcommand\pinkref[1]{{\hypersetup{linkbordercolor=newpink}\autoref{#1}}}

\newtoggle{InString}{}% Keep track of if we are within a string
\togglefalse{InString}% Assume not initally in string

\newcommand*{\ColorIfNotInString}[1]{\iftoggle{InString}{#1}{\color{violet}#1}}%
\newcommand*{\ProcessQuote}[1]{#1\iftoggle{InString}{\global\togglefalse{InString}}{\global\toggletrue{InString}}}%

\lstset{numbers=left,numberblanklines=false,escapeinside=@@, keywordstyle = [2]{\scenvvarstyle}, otherkeywords={arg,this}, morekeywords=[2]{arg,this}}
\let\origthelstnumber\thelstnumber
\makeatletter
\newcommand*\Suppressnumber{
  \lst@AddToHook{OnNewLine}{
    \let\thelstnumber\relax
     \advance\c@lstnumber-\@ne\relax
    }
}

\newcommand*\Reactivatenumber[1]{%
  \setcounter{lstnumber}{\numexpr#1-1\relax}
  \lst@AddToHook{OnNewLine}{%
   \let\thelstnumber\origthelstnumber%
   \refstepcounter{lstnumber}
  }%
}

\makeatother

\begin{document}

\begin{titlepage}
	\begin{center}
	\includegraphics[width=0.2\textwidth]{images/logo.png}\\[0.3cm] 
	\textsc{\normalsize Politecnico di Milano}\\[4cm]
	\line(1,0){300}\\
	[0.25in]
	\huge{\bfseries Euclidean Algorithm for auto-generative patterns in a Supercollider application}\\
	[2mm]
	\line(1,0){200}\\
	[1cm]
	\textsc{\LARGE Project of Computer Music course}\\
	[0.1cm]
	{\large prof. Antonacci Fabio}\\
	[4cm]
	\end{center}
	\begin{flushright}
	\textsc{\large  Students:\\
	de Bari Mauro Giuseppe Matr.899371 \\
	Albertini Davide Matr.883347\\
	January 24, 2019 \\}
	\end{flushright}
\end{titlepage}

\section*{Abstract}
This paper will describe the implementation of a drum-machine Supercollider app, based on Euclidean Rhythms.\\
The \blueref{sec:intro} introduces general concepts about Euclidean Rhythms.\\
The \blueref{sec:algorithm} concerns the historical development of the Euclidean Algorithm, from its first ancient origin to nowday developments for the musical world.\\
The \blueref{sec:app} focuses about the Supercollider application code, describing the code and the algorithms inside the application. For the sake of simplicity, it is divided in three subsection. The first is about the sequencing of instruments, the second is about the humanization of played sequences and the third describes the Graphical User Interface.\\
In the \blueref{sec:concl} informations about the installation and configuration of the app are given. Morover some possible improvements are detailed.
\section{Introduction}
\label{sec:intro}
\subfile{intro}
\section{The algorithm}
\label{sec:algorithm}
\subfile{algorithm}
\section{The application}
\label{sec:app}
In the next three chapter the structure of the application will be explained, with references to Object Oriented Programming. The first section exposes the passage from the creation to the reproduction of a sequence, the second one defines its humanization and the third concludes with an overview on the Graphical User Interface.
\subsection{Sequencing}
\subfile{sequencing}
\subsection{Humanizer}
\subfile{synthdefs}
\subsection{Graphical User Interface}
\subfile{devel}
\section{Conclusions}
\label{sec:concl}
\subfile{conc}
%\nocite{*}
\subfile{bib}

\end{document}